\scanpage{39}
\deleted{\unsure{Rappelons la caractérisation}}

Notations $(\lambda_i(u))$ si $u\in L_0(E)$ : c'est la suite des valeurs \inserted{\deleted{\illegible{3}}} propres de $u$, rangées par ordre de modules non \unsure{décroissants|croissants}, complétée par des $0$ si cette suite est finie.
Si la dim de $E$ est finie, \inserted{soit $n$}, \inserted{il y a lieu d'attacher \deleted{\unsure{et les}} \unsure{l'on range les} \unsure{finies} \unsure{en gardant même} \unsure{du point de vue du rang} $n+1$} \unsure{on joint} \deleted{seulement} à $\lambda_n(u)$ \unsure{les} \unsure{premiers} \unsure{termes} \unsure{nuls}.
\margin{?\unsure{!}}
De plus \deleted{$\rho_n(u)=$}, pour $u\in L_0(E,F)$,
\[
\boxed{\rho_n(u)=\lambda_n(|u|)}
\]
donc $\rho_n(u)$ est une suite décroissante de nbs positifs.

\deleted{$\unsure{\rho_n(u)}\leq\|u\|$}\qquad\deleted{$\unsure{\lambda_n(u)}$}
\[
\boxed{|\lambda_1(u)|\leq\rho_1(u)=\|u\|}\qquad
\boxed{\rho_n(u)=\rho_n(u^*)}
\]
Si $h$ est hermitien positif, on sait que
\[
\boxed{\rho_n(h)=\inf_{E_{n-1}}\deleted{\sup_{\unsure{x\in F}}\unsure{\|h_F\|}}\left(\sup_{\substack{x\in E_{n-1}\\\|x\|\leq1}}(hx,x)\right)^{\unsure{1/2}}}
\]
où \deleted{$\unsure{h_F}$} $E_{n-1}$ parcourt \deleted{\unsure{les espaces} $F:E_{\unsure{n}}$ \unsure{la forme}} les sous-espaces \unsure{vectoriels} fermés de codimension $n-1$.
\deleted{\unsure{contribution} $(hx,y)$}

On a aussi
\[
\rho_n(h^{\unsure{*|2}})^2=\rho_n(h^2)=\inf_{E_n}\sup_{\substack{x\in E_n\\\|x\|\leq1}}\|hx\|^2,
\]
\uline{d'où}
\deleted{$\rho_n(u)=\inf\sup_{x\in E_n}hx$}

\deleted{d'où $\rho_n(u)^2=\rho_n(u^*u)=\inf_{\unsure{E_n}}\sup_{x\in E_n}\|u^*ux\|$}
\[
\rho_n(h)=\inf_{E_n}\sup_{\substack{x\in E_{n-1}\\\|x\|\leq1}}\|hx\|\deleted{\unsure{{}^2}}\qquad(h\geq0)
\]
\deleted{\illegible{1}}
D'où, en prenant $h=|u|$,
\[
\boxed{\rho_n(u)=\inf_{E_n}\sup_{\substack{x\in E_{n-1}\\\|x\|\leq1}}\|ux\|}.
\]
On \deleted{\unsure{en conclut}} \inserted{conclut} aussitôt :
\[
\boxed{\begin{aligned}
\rho_n(uv)&\leq\|u\|\rho_n(v)\\
\rho_n(uv)&\leq\|v\|\rho_n(u)
\end{aligned}}
\]
Et de même : \uline{$\rho_n(h)$ est fonction croissante de l'hermitien positif $h$}.
\deleted{\uline{Corollaire} Soit $u\in L_0(E,F)$}
\deleted{\unsure{En posant} $\boxed{\unsure{\left(\sum\rho_n(u)^p\right)^{1/p}=\|u\|_p}}$}

\deleted{\unsure{Les deux théorèmes} \illegible{1}}
Les inégalités de \unsure{Ky-Fan|ky-Fan}
\[
\boxed{\begin{aligned}
\rho_{\unsure{m+n+1|m+n-1}}(AB)&\leq\rho_{m+1}(A)\rho_{n+1}(B)\\
\rho_{\unsure{m+n+1|m+n-1}}(A+B)&\unsure{<|\leq}\rho_{m+1}(A)+\rho_{n+1}(B)
\end{aligned}}
\]
\margin{?}
\mspage{40}
